\newpage
\Article{54}
Let $h$ be a primitive root for the modulus $m$, of norm $=p$. The terms of the progression
\[
1,h,hh,h^3\dotsfour h^{p-2}
\]
will be incongruent to one another. It follows easily that every integer not divisible by the modulus must be congruent to one of them: this sequence gives a complete system of incongruent residues with zero excluded. The exponent of the power to which a given number is congruent may be called its \emph{index}, while $h$ is regarded as the \emph{base}. Here are some examples, in which we have placed the absolutely least residue beside each index.

\begin{center}\emph{First example.}\end{center}
\[
m=5+4i,\quad p=41,\quad h=1+2i
\]
\begingroup\small
\[
\setlength{\arraycolsep}{2pt}
\begin{array}{r|r||r|r||r|r||r|r||r|r}
\text{Ind.}&\text{Residue}&\text{Ind.}&\text{Residue}&\text{Ind.}&\text{Residue}&\text{Ind.}&\text{Residue}&\text{Ind.}&\text{Residue}\\\hline
0&+1&8&-4&16&-2+2i&24&+2i&32&+1+i\\
1&+1+2i&9&-3+i&17&-1+2i&25&-3i&33&+1+3i\\
2&+1-i&10&-i&18&+4i&26&+2+2i&34&+2\\
3&+3+i&11&+2-i&19&+1+3i&27&+2+i&35&-3\\
4&-2i&12&-1-i&20&-1&28&+4&36&+2-2i\\
5&+3i&13&+1-3i&21&-1-2i&29&+3-i&37&+1-2i\\
6&-2-2i&14&-2&22&-1+i&30&+i&38&-4i\\
7&-2-i&15&+3&23&-3-i&31&-2+i&39&-1-3i
\end{array}
\]
\endgroup

\SourceLeaf{135}%
\begin{center}\emph{Second example.}\end{center}
\[
m=7,\quad p=49,\quad h=1+2i
\]
\begingroup\small
\[
\setlength{\arraycolsep}{2pt}
\begin{array}{r|r||r|r||r|r||r|r||r|r}
\text{Ind.}&\text{Residue}&\text{Ind.}&\text{Residue}&\text{Ind.}&\text{Residue}&\text{Ind.}&\text{Residue}&\text{Ind.}&\text{Residue}\\\hline
0&+1&10&-1-i&20&+2i&30&+2-2i&40&+3\\
1&+1+2i&11&+1-3i&21&+3+2i&31&-1+2i&41&+3-i\\
2&-3-3i&12&-i&22&-1+i&32&+2&42&-2-2i\\
3&+3-2i&13&+2-i&23&-3-i&33&+2-3i&43&+2+i\\
4&-3i&14&-3+3i&24&-1&34&+1+i&44&-2i\\
5&-1-3i&15&-2-3i&25&-1-2i&35&-1+3i&45&-3-2i\\
6&-2+2i&16&-3&26&+3+3i&36&+i&46&+1-i\\
7&+1-2i&17&-3+i&27&-3+2i&37&-2+i&47&+3+i\\
8&-2&18&+2+2i&28&+3i&38&+3-3i&&\\
9&-2+3i&19&-2-i&29&+1+3i&39&+2+3i&&
\end{array}
\]
\endgroup

\Article{55}
We add some observations about primitive roots and the calculation of indices, omitting their proofs because they are easy.

I. In a given system, indices congruent modulo $p-1$ correspond to residues congruent modulo $m$, and conversely.

II. Residues corresponding to indices relatively prime to $p-1$ are themselves primitive roots, and conversely.

III. Suppose that, with a primitive root $h$ taken as the base, the index of another primitive root $h'$ is $t$, and, conversely, $t'$ is the index of $h$ when $h'$ is taken as the base. Then $tt'\equiv1\pmod{p-1}$. Under the same assumptions, if the indices of some other number in these two systems are respectively $u,u'$, then $tu'\equiv u,t'u\equiv u'\pmod{p-1}$.

IV. Excluding the numbers $1,1+i$ and their three associates each, as being too trivial, from the moduli we consider leaves the prime numbers that we placed third and fourth in Article~34. The norms of the latter are real prime numbers of the form $4n+1$; those of the former, however, are squares of odd real prime numbers. Thus, in either case, $p-1$ is divisible by $4$.

V. Denote the index of $-1$ by $u$. Then $2u\equiv0\pmod{p-1}$, so either $u\equiv0$ or $u\equiv\frac12(p-1)$. But since the index $0$ corresponds to the residue $+1$, the index of $-1$ must be $\frac12(p-1)$.

\begingroup\rightskip=0pt plus 1fil\relax
VI. Similarly, denoting by $u$ the index of $i$, we have \mbox{$2u\equiv\frac12(p-1)\pmod{p-1}$}, and therefore either $u\equiv\frac14(p-1)$ or $u\equiv\frac34(p-1)$. Here, however, the ambiguity depends on the choice of primitive root. Specifically, if, with the primitive root $h$ taken as the base,
\SourceLeaf{136}%
the index of $i$ is $\frac14(p-1)$, the index becomes $\frac34(p-1)$ when $h^\mu$ is taken as the base, where $\mu$ denotes a positive integer of the form $4n+3$ relatively prime to $p-1$, for example the number $p-2$ itself; and conversely. Thus one half of the primitive roots gives $i$ the index $\frac14(p-1)$, and the other half the index $\frac34(p-1)$. Evidently, for the former bases $-i$ has the index $\frac34(p-1)$, and for the latter the index $\frac14(p-1)$.
\par\endgroup
VII. Whenever the modulus is a positive real prime number of the form $4n+3$, say $=q$, so that $p=qq$, the indices of all real numbers are divisible by $q+1$. For, if $t$ denotes the index of the real number $k$, then, because $k^{q-1}\equiv\pmod q$, we have $(q-1)t\equiv0\pmod{qq-1}$, and consequently $\frac{t}{q+1}$ is an integer. Similarly, the indices of purely imaginary numbers such as $ki$ are divisible by $\frac12(q+1)$. Thus primitive roots for such moduli must evidently be sought only among mixed numbers.

VIII. For a modulus $m$ that is a mixed complex prime number, however, whose norm $p$ is consequently a real prime number of the form $4n+1$, any primitive roots may also be chosen among real numbers, among which a complete system of incongruent residues can indeed be exhibited (Article~40). Every real number that is a primitive root for the complex modulus $m$ is evidently also, in the arithmetic of real numbers, a primitive root for the modulus $p$, and conversely.

